\documentclass[tikz,border=3mm,multi=tikzpicture]{standalone}
\usepackage{eleve-fisica}

\begin{document}

% FIGURA 1 — fig-01-linhas-de-campo-do-ima
\begin{tikzpicture}[fisica figura, x=1cm, y=1cm]
  \path[use as bounding box] (-0.2,-0.2) rectangle (11.0,8.6);
  \node[fisica corpo, minimum width=4.2cm, minimum height=1.45cm] (I)
    at (5.25,4.0) {};
  \draw[elevecinza, line width=1.2pt] (5.25,3.28) -- (5.25,4.72);
  \node[font=\Large\bfseries, text=eleveazul] at (4.2,4.0) {S};
  \node[font=\Large\bfseries, text=elevelaranja] at (6.3,4.0) {N};
  \draw[fisica campo] (6.95,4.55)
    .. controls (9.2,6.25) and (1.3,6.25) .. (3.55,4.55);
  \draw[fisica campo] (6.95,4.2)
    .. controls (8.65,5.15) and (1.85,5.15) .. (3.55,4.2);
  \draw[fisica campo] (6.95,3.45)
    .. controls (9.2,1.75) and (1.3,1.75) .. (3.55,3.45);
  \draw[fisica campo] (6.95,3.8)
    .. controls (8.65,2.85) and (1.85,2.85) .. (3.55,3.8);
  \draw[fisica campo] (3.55,4.0) -- (6.95,4.0);
  \node[fisica rotulo] at (5.25,7.55) {fora: N $\to$ S};
  \node[fisica rotulo] at (5.25,0.65) {dentro: S $\to$ N};
\end{tikzpicture}

% FIGURA 2 — fig-02-terra-como-dipolo-inclinado
\begin{tikzpicture}[fisica figura, x=1cm, y=1cm]
  \path[use as bounding box] (-0.2,-0.2) rectangle (10.2,9.7);
  \coordinate (O) at (4.8,4.55);
  \shade[ball color=eleveazulclaro] (O) circle (2.5);
  \draw[eleveazul, line width=1.3pt] (O) circle (2.5);
  \draw[elevecinza, line width=1.4pt] ($(O)+(0,-3.8)$) -- ($(O)+(0,3.8)$);
  \draw[elevelaranja, line width=1.7pt] ($(O)+(-0.73,-3.73)$) -- ($(O)+(0.73,3.73)$);
  \node[fisica rotulo, anchor=west] at (5.05,8.65) {eixo geográfico};
  \node[fisica medida, anchor=west] at (5.95,7.65) {eixo magnético};
  \draw[elevelaranja, line width=1.1pt] ($(O)+(0,2.95)$)
    arc[start angle=90,end angle=79,radius=2.95];
  \node[fisica medida] at (5.55,7.0) {$\approx11^\circ$};
  \draw[fisica campo] ($(O)+(0.65,2.4)$)
    .. controls (8.7,7.0) and (8.7,2.1) .. ($(O)+(-0.65,-2.4)$);
  \draw[fisica campo] ($(O)+(-0.65,-2.4)$)
    .. controls (0.9,2.1) and (0.9,7.0) .. ($(O)+(0.65,2.4)$);
  \node[font=\large\bfseries, text=elevecinza] at (4.8,0.25)
    {modelo de dipolo};
\end{tikzpicture}

% FIGURA 3 — fig-03-declinacao-e-inclinacao
\begin{tikzpicture}[fisica figura, x=1cm, y=1cm]
  \path[use as bounding box] (-0.2,-0.2) rectangle (10.8,10.3);
  \node[font=\Large\bfseries, text=eleveazul] at (5.15,9.75)
    {vista horizontal: declinação};
  \coordinate (D) at (5.15,6.85);
  \draw[fisica eixo] (D) -- ++(0,2.1)
    node[above, fisica rotulo] {norte geográfico};
  \draw[fisica campo, ->] (D) -- ++(-1.35,1.7)
    node[fisica rotulo, above left] {$\vec B_h$};
  \draw[elevelaranja, line width=1.1pt] ($(D)+(0,1.15)$)
    arc[start angle=90,end angle=128,radius=1.15];
  \node[fisica medida] at (4.45,7.95) {$D$};

  \node[font=\Large\bfseries, text=eleveazul] at (5.15,4.75)
    {vista vertical: inclinação};
  \coordinate (I) at (5.15,1.8);
  \draw[fisica eixo] (I) -- ++(2.55,0)
    node[right, fisica rotulo] {horizontal};
  \draw[fisica campo, ->] (I) -- ++(1.85,-1.35)
    node[fisica rotulo, below right] {$\vec B$};
  \draw[elevelaranja, line width=1.1pt] ($(I)+(1.15,0)$)
    arc[start angle=0,end angle=-36,radius=1.15];
  \node[fisica medida] at (6.45,1.25) {$I$};
\end{tikzpicture}

% FIGURA 4 — fig-04-campo-ao-redor-do-fio
\begin{tikzpicture}[fisica figura, x=1cm, y=1cm]
  \path[use as bounding box] (-0.2,-0.2) rectangle (9.8,9.4);
  \coordinate (O) at (4.65,4.45);
  \node[font=\Huge\bfseries, text=elevelaranja] at (O) {$\odot$};
  \node[fisica medida, below=9pt] at (O) {$I$ saindo do plano};
  \foreach \r in {1.3,2.35,3.4}
    \draw[fisica campo] ($(O)+(0:\r)$) arc[start angle=0,end angle=360,radius=\r];
  \draw[fisica movimento] (7.8,7.4)
    arc[start angle=35,end angle=145,radius=1.0];
  \node[fisica rotulo] at (4.65,8.75) {campo anti-horário};
  \node[fisica rotulo] at (8.35,4.45) {$\vec B$};
\end{tikzpicture}

% FIGURA 5 — fig-05-campo-da-espira
\begin{tikzpicture}[fisica figura, x=1cm, y=1cm]
  \path[use as bounding box] (-0.2,-0.2) rectangle (10.6,8.9);
  \coordinate (O) at (5.0,3.95);
  \draw[elevecinza, line width=2pt] (O) ellipse (3.6 and 1.8);
  \draw[fisica movimento] (2.3,2.75)
    arc[start angle=215,end angle=320,x radius=3.6,y radius=1.8]
    node[fisica rotulo, right] {$I$};
  \fill[eleveazul] (O) circle (3pt);
  \draw[fisica campo, ->] (O) -- ++(0,3.7)
    node[fisica rotulo, above] {$\vec B$};
  \draw[fisica auxiliar] (O) -- ++(0,-2.8);
  \node[fisica rotulo] at (7.8,1.1) {plano da espira};
  \draw[fisica auxiliar, ->] (7.25,1.4) -- (6.55,2.45);
  \node[font=\large\bfseries, text=elevecinza] at (5.0,8.25)
    {$\vec B$ perpendicular no centro};
\end{tikzpicture}

% FIGURA 6 — fig-06-campo-do-solenoide
\begin{tikzpicture}[fisica figura, x=1cm, y=1cm]
  \path[use as bounding box] (-0.2,-0.2) rectangle (11.2,8.8);
  \foreach \x in {2.0,2.65,...,8.5}
    \draw[elevecinza, line width=1.2pt] (\x,2.1) arc[start angle=270,end angle=90,x radius=.45,y radius=2.0];
  \foreach \y in {2.8,4.05,5.3}
    \draw[fisica campo] (1.35,\y) -- (9.15,\y);
  \draw[fisica campo] (9.15,5.8)
    .. controls (10.6,7.6) and (0.3,7.6) .. (1.35,5.8);
  \draw[fisica campo] (9.15,2.0)
    .. controls (10.6,0.2) and (0.3,0.2) .. (1.35,2.0);
  \node[fisica rotulo] at (5.25,6.65) {campo interno quase uniforme};
  \node[fisica rotulo] at (5.25,0.15) {linhas externas retornam};
  \node[fisica medida] at (0.75,4.05) {S};
  \node[fisica medida] at (9.75,4.05) {N};
\end{tikzpicture}

% FIGURA 7 — fig-07-experiencia-de-oersted
\begin{tikzpicture}[fisica figura, x=1cm, y=1cm]
  \path[use as bounding box] (-0.2,-0.2) rectangle (10.8,9.1);
  \draw[elevecinza, line width=1.7pt] (1.0,7.25) -- (9.65,7.25);
  \draw[fisica movimento] (2.1,7.25) -- (5.1,7.25)
    node[fisica rotulo, above] {$I$};
  \draw[elevecinza, line width=1.7pt] (1.0,7.25) -- (1.0,1.15)
    -- (3.45,1.15);
  \draw[elevecinza, line width=1.7pt] (6.95,1.15) -- (9.65,1.15) -- (9.65,7.25);
  \draw[elevecinza, line width=1.3pt] (3.45,0.65) rectangle (6.95,1.65);
  \node[font=\large\bfseries, text=elevecinza] at (5.2,1.15) {pilha};
  \draw[eleveazul, line width=1.2pt] (5.2,3.0) circle (1.15);
  \draw[elevelaranja, line width=2.2pt] (4.4,2.35) -- (6.0,3.65);
  \fill[elevelaranja] (6.0,3.65) -- ++(-0.45,-0.05) -- ++(0.15,-0.4) -- cycle;
  \node[fisica rotulo] at (5.2,4.6) {bússola desviada};
  \draw[fisica campo] (7.55,5.55) arc[start angle=20,end angle=160,radius=2.45];
  \node[font=\large\bfseries, text=elevecinza] at (5.2,8.45)
    {a corrente cria campo circular};
\end{tikzpicture}

\end{document}
